\chapter{Novidade no espaço latente}
\label{cap:mahalanobis}

\begin{objetivos}
Ao final deste capítulo, você será capaz de:
\begin{itemize}[leftmargin=*,itemsep=0.2em]
  \item aplicar a distância de Mahalanobis como escore de novidade e calibrar seu limiar
        por percentil;
  \item comparar espaços de representação --- bruto, latente de autoencoder e penúltima
        camada de classificador --- por índices internos de qualidade;
  \item explicar por que uma representação treinada para discriminar supera uma
        treinada para reconstruir, na tarefa de detectar o desconhecido;
  \item interpretar razões de separação entre classes conhecidas e desconhecidas;
  \item reconhecer as hipóteses estatísticas embutidas no método e as situações em que
        elas falham.
\end{itemize}
\end{objetivos}

\section{Retomando um resultado incômodo}

O \cref{cap:mundoaberto} apresentou um achado que ficou sem explicação: autoencoder,
VAE e DCN produziram acurácias de $0{,}730$, $0{,}727$ e $0{,}725$. Três arquiteturas de
sofisticação crescente, desempenho indistinguível.

A interpretação oferecida na ocasião foi que o gargalo estava em outro lugar. Este
capítulo mostra que a interpretação estava certa e localiza o gargalo com precisão: o
problema não é \emph{qual} autoencoder, mas o fato de usar um autoencoder.

\begin{destaque}[A hipótese deste capítulo]
Um autoencoder é treinado para minimizar erro de reconstrução. Nada em seu objetivo o
incentiva a \emph{separar} classes; ao contrário, ele gasta capacidade representando
fielmente aquilo que é comum a todas elas.

Um classificador multiclasse é treinado justamente para separar. A representação
implícita em sua penúltima camada --- o vetor que alimenta a camada de saída --- é,
por construção, um espaço em que as classes são linearmente separáveis.

A hipótese é que \textbf{esse espaço, construído para discriminar entre classes
conhecidas, também é melhor para reconhecer o desconhecido} --- mesmo que nada em seu
treinamento tenha visado a isso.
\end{destaque}

\section{A distância de Mahalanobis como escore}

\subsection{Definição}

Dado um conjunto de referência com média $\bm{\mu}$ e covariância $\mat{\Sigma}$, a
distância de Mahalanobis (\cref{eq:mahalanobis}) de um ponto $\vect{z}$ é
\[
  d_M(\vect{z}) = \sqrt{(\vect{z} - \bm{\mu})^\top \mat{\Sigma}^{-1} (\vect{z} - \bm{\mu})}.
\]

Sua virtude sobre a distância euclidiana é ponderar cada direção pela variabilidade
observada. Em um espaço latente onde algumas dimensões variam muito naturalmente e
outras quase nada, essa correção é o que separa um escore útil de um escore dominado
pelas direções de maior escala.

\subsection{Calibração do limiar}

O limiar é fixado no \textbf{percentil 95} da distribuição de $d_M$ sobre os dados
conhecidos. A escolha é operacional, não estatística: aceita-se, por construção, uma
taxa de falso alarme de aproximadamente $5\,\%$ em troca de sensibilidade.

\begin{destaque}[Por que percentil e não $\chi^2$]
Sob normalidade multivariada, $d_M^2$ segue distribuição $\chi^2$ com graus de liberdade
iguais à dimensão, e o limiar poderia ser lido diretamente de uma tabela.

Os dados não são normais. Calibrar por percentil empírico dispensa a hipótese e produz a
taxa de falso alarme desejada \emph{por construção}, qualquer que seja a distribuição
real. É um exemplo geral: quando a hipótese paramétrica é duvidosa e há dados
suficientes, o quantil empírico é preferível ao valor teórico.
\end{destaque}

\section{Três espaços em competição}

\begin{enumerate}[leftmargin=*,itemsep=0.4em]
  \item \textbf{Espaço bruto}, com $200 \times 4 = 800$ dimensões. Nenhuma transformação.
  \item \textbf{Latente de autoencoder LSTM}, com $128$ dimensões. Codificador:
        entrada $\to$ LSTM(256) $\to$ LSTM(128) $\to$ Densa(128). Decodificador
        simétrico. Treinado para reconstruir.
  \item \textbf{Penúltima camada de classificador multiclasse}, com $64$ dimensões.
        Treinada para discriminar entre as classes conhecidas.
\end{enumerate}

\begin{figure}[htbp]
\centering
\includegraphics[width=\textwidth]{binary_classifiers.png}
\caption{Cabeça classificadora aplicada sobre a representação latente: duas camadas
densas com ativação ReLU seguidas da camada de saída. É a mesma cabeça para os três
espaços comparados, o que isola o efeito da representação.}
\label{fig:mahal-cabeca}
\end{figure}

O controle experimental importa: a arquitetura da cabeça classificadora é idêntica nos
três casos. Qualquer diferença de desempenho é atribuível à representação, não ao
classificador.

\section{Resultados}

\subsection{Detecção de novidade}

A \cref{tab:mahal-deteccao} apresenta o resultado central. Cada linha corresponde a uma
classe tratada como \emph{desconhecida} --- removida do treinamento --- e reporta a
fração de suas amostras corretamente identificadas como novidade por três métodos.

\begin{table}[htbp]
\centering
\caption{Acurácia de detecção de novidade (\%) por classe tratada como desconhecida,
sob três abordagens.}
\label{tab:mahal-deteccao}
\small
\begin{tabular}{lccc}
\toprule
\textbf{Classe tratada como desconhecida} & \textbf{Class.\ bin.} &
\textbf{Mahal.\ bruto} & \textbf{Mahal.\ latente} \\
\midrule
1 --- Aumento de BSW           & $22{,}8$ & $41{,}5$ & $\bm{99{,}0}$ \\
2 --- Fechamento de DHSV       & $76{,}9$ & $77{,}1$ & $\bm{96{,}5}$ \\
3 --- Golfadas severas         & $48{,}1$ & $61{,}0$ & $\bm{99{,}0}$ \\
4 --- Instabilidade de fluxo   & $44{,}7$ & $57{,}8$ & $\bm{96{,}0}$ \\
5 --- Perda rápida de produtiv.& $\bm{0{,}1}$ & $62{,}3$ & $67{,}5$ \\
6 --- Restrição rápida no PCK  & $47{,}6$ & $47{,}6$ & $\bm{99{,}0}$ \\
7 --- Incrustação no PCK       & $36{,}5$ & $36{,}7$ & $52{,}3$ \\
8 --- Hidrato na produção      & $50{,}5$ & $91{,}1$ & $\bm{99{,}0}$ \\
9 --- Hidrato no serviço       & $51{,}9$ & $88{,}7$ & $\bm{99{,}0}$ \\
Falha de ICV                   & $37{,}0$ & $67{,}5$ & $\bm{98{,}2}$ \\
\bottomrule
\end{tabular}
\end{table}

\begin{resultadochave}
Aplicar a distância de Mahalanobis no espaço latente do classificador multiclasse eleva
a detecção de novidade a $99\,\%$ em cinco das dez classes, e a mais de $96\,\%$ em sete
delas. Sobre a anomalia de ICV --- a única genuinamente inédita --- o salto é de
$52\,\%$ (\cref{cap:mundoaberto}) para $\bm{98{,}2\,\%}$.
\end{resultadochave}

\subsection{Três observações}

\paragraph{A progressão é sistemática.} Em nove das dez linhas, o Mahalanobis no espaço
bruto supera os classificadores binários, e o Mahalanobis latente supera o bruto. Não é
um ganho concentrado em casos favoráveis: é uma melhoria uniforme, o que é a assinatura
de uma mudança correta de método e não de ajuste ao conjunto de teste.

\paragraph{A classe 5 e o valor de $0{,}1\,\%$.} Um detector que identifica um décimo de
um por cento das novidades de uma classe está, para todos os efeitos, cego a ela. Ainda
assim, o Mahalanobis latente eleva o número a $67{,}5\,\%$ --- o pior resultado da
coluna, e um ganho de seiscentas vezes. A classe 5 continua sendo o ponto fraco do
sistema em todos os capítulos deste livro.

\paragraph{A classe 7 e o limite físico.} A incrustação no PCK atinge apenas $52{,}3\,\%$
mesmo no melhor espaço. A explicação é a escala temporal: incrustação se desenvolve ao
longo de semanas, e uma janela de duzentos passos não contém a evidência necessária.
Nenhuma escolha de representação corrige um problema de janela.

\subsection{Quanto o desconhecido se afasta}

A acurácia diz se o limiar foi ultrapassado. A \cref{tab:mahal-separacao} diz por quanto.

\begin{table}[htbp]
\centering
\caption{Razão de separação: distância de Mahalanobis média das amostras desconhecidas
dividida pela das conhecidas, nos espaços bruto e latente.}
\label{tab:mahal-separacao}
\small
\begin{tabular}{lcc}
\toprule
\textbf{Classe desconhecida} & \textbf{Espaço bruto} & \textbf{Espaço latente} \\
\midrule
1 & $2{,}1\times$ & $15{,}8\times$ \\
2 & $1{,}8\times$ & $12{,}4\times$ \\
3 & $2{,}3\times$ & $18{,}2\times$ \\
4 & $2{,}0\times$ & $16{,}1\times$ \\
5 & $1{,}9\times$ & $14{,}7\times$ \\
6 & $2{,}4\times$ & $13{,}9\times$ \\
8 & $<1{,}0\times$ & $\bm{186{,}5\times}$ \\
9 & $2{,}2\times$ & $17{,}3\times$ \\
\bottomrule
\end{tabular}
\end{table}

\begin{destaque}[A linha da classe 8]
No espaço bruto, a razão de separação da classe 8 é \emph{menor que um}: as amostras de
hidrato estão, em média, \textbf{mais próximas} do centro dos dados conhecidos do que as
próprias amostras conhecidas. Qualquer detector de novidade baseado em distância
falharia --- não por má calibração, mas porque no espaço bruto a informação simplesmente
não está disponível de forma acessível a uma métrica de distância.

No espaço latente, a mesma classe atinge $186{,}5\times$. É o maior valor da tabela por
uma ordem de grandeza.

A inversão é a demonstração mais forte da tese deste capítulo. Nenhum ajuste de limiar,
nenhuma escolha de escore, nenhuma sofisticação algorítmica recuperaria o desempenho no
espaço bruto. \textbf{A representação não é um detalhe de pré-processamento; ela
determina o que é detectável.}
\end{destaque}

Vale ainda notar a uniformidade das demais linhas: entre $12{,}4\times$ e $18{,}2\times$
no latente, contra $1{,}8\times$ a $2{,}4\times$ no bruto. O ganho é de aproximadamente
uma ordem de grandeza, de forma consistente.

\subsection{Por que o latente é melhor: a evidência direta}

Até aqui o argumento foi indireto --- o desempenho é melhor, logo o espaço é melhor. A
\cref{tab:mahal-qualidade} mede a qualidade dos espaços diretamente, pelos índices
internos do \cref{cap:novidades}.

\begin{table}[htbp]
\centering
\caption{Qualidade do espaço de representação segundo três índices internos. Para
silhueta e Calinski-Harabasz, maior é melhor; para Davies-Bouldin, menor é melhor.}
\label{tab:mahal-qualidade}
\small
\begin{tabular}{lcc}
\toprule
\textbf{Índice} & \textbf{Latente do autoencoder} & \textbf{Latente do multiclasse} \\
\midrule
Silhueta            & $-0{,}0149$   & $\bm{0{,}2299}$ \\
Calinski-Harabasz   & $687{,}8008$  & $\bm{2818{,}231}$ \\
Davies-Bouldin      & $4{,}5592$    & $\bm{1{,}746}$ \\
\bottomrule
\end{tabular}
\end{table}

\begin{resultadochave}
A silhueta do latente do autoencoder é \textbf{negativa} ($-0{,}0149$): em média, as
amostras estão mais próximas de um agrupamento vizinho do que do seu próprio. O latente
do multiclasse atinge $0{,}2299$. O índice Calinski-Harabasz é quatro vezes maior e o
Davies-Bouldin, dois pontos e meio menor.
\end{resultadochave}

\begin{destaque}[A explicação retroativa do \cref{cap:mundoaberto}]
Silhueta negativa significa que \textbf{não há estrutura de classe} no latente do
autoencoder. Isso explica, de uma vez, três observações anteriores:

por que autoencoder, VAE e DCN produziram desempenho equivalente --- todos otimizavam
reconstrução, e nenhum criava separação;

por que o agrupamento sobre o material rejeitado alcançou apenas $80{,}8\,\%$ --- não há
agrupamentos bem definidos para encontrar;

e por que o refinamento por limiar adaptativo precisou de tanta engenharia --- ele estava
compensando a ausência de estrutura no espaço subjacente.

O esforço do \cref{cap:mundoaberto} foi, em boa medida, gasto tentando extrair estrutura
de um espaço que não a tinha.
\end{destaque}

\subsection{Efeito sobre a classificação}

O ganho não se restringe à detecção de novidade. A \cref{tab:mahal-classificacao} mostra
o desempenho dos classificadores binários --- arquitetura idêntica --- sobre os três
espaços.

\begin{table}[htbp]
\centering
\caption{Desempenho dos classificadores binários em função do espaço de representação.
Média e desvio padrão sobre as classes.}
\label{tab:mahal-classificacao}
\small
\begin{tabular}{lcccc}
\toprule
\textbf{Espaço} & \textbf{ACC} & \textbf{PR} & \textbf{REC} & $\bm{F_1}$ \\
\midrule
Latente do multiclasse & $\bm{0{,}755} \pm 0{,}171$ & $\bm{0{,}732} \pm 0{,}280$ &
                          $\bm{0{,}685} \pm 0{,}384$ & $\bm{0{,}668} \pm 0{,}322$ \\
Latente do autoencoder & $0{,}697 \pm 0{,}154$ & $0{,}697 \pm 0{,}277$ &
                          $0{,}530 \pm 0{,}367$ & $0{,}559 \pm 0{,}318$ \\
Espaço bruto           & $0{,}691 \pm 0{,}162$ & $0{,}642 \pm 0{,}329$ &
                          $0{,}535 \pm 0{,}392$ & $0{,}543 \pm 0{,}349$ \\
\bottomrule
\end{tabular}
\end{table}

A ordem se mantém, mas as diferenças são bem menores --- seis pontos percentuais de
acurácia entre o melhor e o pior espaço, contra dezenas de pontos na detecção de
novidade.

\begin{destaque}[Uma assimetria reveladora]
A qualidade do espaço importa muito mais para \emph{detectar o desconhecido} do que para
\emph{classificar o conhecido}.

A razão é que um classificador suficientemente expressivo pode compensar uma
representação ruim: ele aprende a fronteira onde quer que ela esteja. Um escore de
distância não pode --- ele depende inteiramente da geometria que lhe é dada.

Corolário prático: se o objetivo é classificar bem o que se conhece, invista no
classificador. Se o objetivo é reconhecer o que não se conhece, invista na
representação.
\end{destaque}

Note também os desvios padrão elevados --- revocação de $0{,}685 \pm 0{,}384$. Um desvio
maior que metade da média indica que o desempenho varia enormemente entre classes, o que
as tabelas por classe dos capítulos anteriores já mostraram: excelente para as classes
2, 3 e 9, ruim para as classes 4 e 5.

\section{Balanço da evolução}

\begin{resultadochave}
No conjunto das classes avaliadas, a distância de Mahalanobis no espaço latente do
classificador multiclasse detecta $\bm{95{,}4\,\%}$ das falhas não vistas com
$\bm{4{,}5\,\%}$ de falsos alarmes --- uma melhoria média de $17\,\%$ sobre a melhor
alternativa anterior.
\end{resultadochave}

A \cref{tab:mahal-evolucao} resume o percurso da \cref{parte:evolucao}.

\begin{table}[htbp]
\centering
\caption{Evolução da detecção de novidade sobre a anomalia de ICV ao longo do livro.}
\label{tab:mahal-evolucao}
\small
\begin{tabular}{lll}
\toprule
\textbf{Capítulo} & \textbf{Método} & \textbf{Detecção} \\
\midrule
\cref{cap:dual}         & Erro de reconstrução, limiar $\mu + 3\sigma$ & apenas binária \\
\cref{cap:mundoaberto}  & Ensemble binário no latente do autoencoder   & $52{,}0\,\%$ \\
\cref{cap:mahalanobis}  & Mahalanobis no espaço bruto                  & $67{,}5\,\%$ \\
\cref{cap:mahalanobis}  & Mahalanobis no latente do multiclasse        & $\bm{98{,}2\,\%}$ \\
\bottomrule
\end{tabular}
\end{table}

\begin{destaque}[Uma lição sobre engenharia de sistemas de aprendizado]
A progressão de $52\,\%$ a $98{,}2\,\%$ não veio de um modelo maior, de mais dados ou de
uma arquitetura nova. Veio de duas mudanças conceituais: segmentar em vez de detectar
(\cref{cap:segmentacao}) e representar com um classificador em vez de um autoencoder
(este capítulo).

Ambas custam pouco em implementação. Ambas exigiram identificar corretamente onde estava
o gargalo --- e, no caso da representação, foi um resultado \emph{negativo} do
\cref{cap:mundoaberto} que apontou o caminho. Resultados negativos, quando reportados,
são informação.
\end{destaque}

\section{Limitações}

\begin{enumerate}[leftmargin=*,itemsep=0.4em]

\item \textbf{Hipótese elíptica.} A distância de Mahalanobis supõe distribuição
aproximadamente elíptica, caracterizada por média e covariância. Se a classe conhecida é
multimodal --- e várias são ---, a hipótese é violada e uma mistura de gaussianas ou uma
estimativa por vizinhança seria mais adequada.

\item \textbf{Ausência de tratamento de deriva.} $\bm{\mu}$ e $\mat{\Sigma}$ são
estimados uma vez. Poços mudam ao longo de anos, e o método não prevê atualização.

\item \textbf{Dependência da segmentação.} Todo o desempenho pressupõe segmentos bem
delimitados. Nos poços em que a U-Net vai mal --- B1, com IoU de $56\,\%$ ---, o escore
de novidade é calculado sobre um segmento contaminado.

\item \textbf{Variabilidade entre poços.} Persistente, como os desvios padrão da
\cref{tab:mahal-classificacao} mostram.

\item \textbf{Revocação e antecedência.} SoftED de $0{,}678$ e antecedência de $0{,}222$
permanecem inalteradas por este capítulo, que atua sobre a natureza da decisão e não
sobre o instante em que ela é tomada.

\item \textbf{Classes 5 e 7.} $67{,}5\,\%$ e $52{,}3\,\%$ mesmo no melhor espaço, por
razões que são de rótulo e de escala temporal, respectivamente.

\item \textbf{Ainda não há explicação.} O sistema agora detecta o desconhecido com
$98{,}2\,\%$ de acerto e continua incapaz de dizer uma palavra sobre o que ele é.

\end{enumerate}

\begin{destaque}[O que a \cref{parte:evolucao} deixa para a seguinte]
Ao final desta parte, o sistema sabe \emph{onde} está o evento, sabe \emph{que} ele é
desconhecido, e sabe com alta confiança. A única pergunta do \cref{cap:introducao} que
permanece inteiramente sem resposta é a terceira: \emph{por quê?}
\end{destaque}

\resumocapitulo{%
A distância de Mahalanobis aplicada ao espaço latente extraído da penúltima camada de um
classificador multiclasse detecta $95{,}4\,\%$ das falhas não vistas com $4{,}5\,\%$ de
falsos alarmes, contra $67{,}5\,\%$ no espaço bruto e $52\,\%$ com o ensemble sobre o
latente do autoencoder. Sobre a anomalia de ICV, a detecção salta de $52\,\%$ para
$98{,}2\,\%$. A razão é geométrica e mensurável: a silhueta do latente do autoencoder é
negativa ($-0{,}0149$), indicando ausência de estrutura de classe, enquanto a do
multiclasse é $0{,}2299$, com Calinski-Harabasz quatro vezes maior e Davies-Bouldin dois
pontos e meio menor. A razão de separação entre desconhecidos e conhecidos passa de
cerca de $2\times$ no bruto para $12$ a $18\times$ no latente, e atinge $186{,}5\times$
para a classe 8 --- que no espaço bruto tinha razão \emph{menor que um}. A qualidade da
representação importa muito mais para detectar o desconhecido do que para classificar o
conhecido: seis pontos percentuais de acurácia de classificação contra dezenas de pontos
de detecção de novidade. Isso explica retroativamente por que autoencoder, VAE e DCN
produziram desempenho indistinguível no \cref{cap:mundoaberto}. Permanecem em aberto as
classes 5 e 7, a antecedência de detecção e --- sobretudo --- a explicação.}

\begin{exercicios}
\item Demonstre que a distância de Mahalanobis é invariante a transformações lineares
      invertíveis do espaço, e que a distância euclidiana não é. Relacione esse fato à
      superioridade observada da primeira em espaços latentes com escalas heterogêneas.

\item A razão de separação da classe 8 é menor que $1{,}0$ no espaço bruto. Construa um
      exemplo bidimensional em que amostras de uma classe ausente do treinamento fiquem,
      em média, mais próximas do centroide dos dados de treinamento do que os próprios
      dados de treinamento. Que geometria produz esse efeito?

\item Implemente os três espaços e reproduza a \cref{tab:mahal-qualidade}. Em seguida,
      calcule a silhueta do latente de um classificador multiclasse treinado com
      diferentes números de classes ($3$, $5$, $7$, $9$). A qualidade do espaço melhora
      monotonicamente com o número de classes?

\item A hipótese elíptica é violada quando a classe conhecida é multimodal. Implemente
      uma alternativa baseada em mistura de gaussianas e compare com o Mahalanobis
      simples sobre as classes do 3W. Em quais classes o ganho é maior? Isso corresponde
      à sua expectativa sobre quais classes são multimodais?

\item O limiar é o percentil 95 dos dados conhecidos. Trace a curva de detecção de
      novidade \emph{versus} taxa de falso alarme variando o percentil de 80 a 99{,}9,
      para o espaço bruto e para o latente. Compare as áreas sob as curvas.

\item Argumenta-se que a classe 7 falha por razão de escala temporal e não de
      representação. Projete um experimento que teste essa hipótese --- por exemplo,
      repetindo a avaliação com janelas progressivamente maiores --- e preveja o
      resultado esperado sob a hipótese e sob sua negação.

\item Proponha um mecanismo de atualização de $\bm{\mu}$ e $\mat{\Sigma}$ que acompanhe
      a deriva lenta do poço sem absorver eventos anômalos. Discuta o compromisso entre
      taxa de adaptação e risco de o modelo aprender a considerar normal uma degradação
      progressiva.
\end{exercicios}

\begin{leituras}
\item \citet{Lopes2026OTC} --- o artigo que originou este capítulo.
\item \citet{Mahalanobis1936} --- o trabalho original que define a distância.
\item \citet{Rousseeuw1987}, \citet{Calinski1974} --- os índices internos usados na
      comparação de espaços.
\item \citet{Pimentel2014review} --- revisão de detecção de novidade, útil para situar o
      Mahalanobis entre as alternativas.
\item \citet{maaten2008visualizing} --- para a visualização dos espaços latentes, com as
      ressalvas já discutidas.
\end{leituras}
